\documentclass[11pt,a4paper]{article}
\usepackage[T1]{fontenc}
\usepackage{lmodern,amsmath,amssymb,amsthm,microtype}
\usepackage[margin=27mm]{geometry}
\usepackage[colorlinks=true,linkcolor=blue,citecolor=blue,urlcolor=blue]{hyperref}
\newtheorem{theorem}{Theorem}[section]
\newtheorem{corollary}[theorem]{Corollary}
\newtheorem{proposition}[theorem]{Proposition}
\newcommand{\Z}{\mathbb Z}\newcommand{\Q}{\mathbb Q}
\setlength{\emergencystretch}{2em}
\title{A reduction obstruction to adjacent contacts on rank-stable elliptic curves}
\author{}\date{17 November 2026}
\begin{document}\maketitle
\begin{abstract}
A non-torsion elliptic point that descends to a smaller number field
after a fixed multiplication cannot realize adjacent index contacts
at two places above the same prime, except at finitely many primes.
The assertion includes torsion translates and any fixed finite set
of nonzero index shifts. Consequently a rank-stable elliptic curve
defined over the base field cannot be substituted directly into the
conjugate adjacent-index detector used in the rational construction.
This is an obstruction to that local completeness mechanism, not to
Diophantine definability of rings of integers.
\end{abstract}

\section{Two reductions above one prime}
Let $L/K$ be a finite extension of number fields, let $E/K$ be an
elliptic curve, and let $P\in E(L)$ have infinite order. Suppose
that $[t_0]P\in E(K)$ for a positive integer $t_0$.
Choose a positive multiple $t$ of $t_0$ annihilating
$E(L)_{\rm tors}$, and put $Q=[t]P\in E(K)$. The point $Q$ is
non-torsion. For a prime of good reduction $\mathfrak p$ of $K$,
write $\widetilde Q_{\mathfrak p}$ for the reduction of $Q$.
For a place $\mathfrak P$ of $L$ above $\mathfrak p$, reduction
of any $K$-rational point agrees with extension of its reduction
from the residue field of $\mathfrak p$ to that of $\mathfrak P$.

\begin{theorem}\label{thm:reduction}
Let $c\in\Z\setminus\{0\}$. Suppose that for some integer $h$,
some torsion points $T_1,T_2\in E(L)$, and two places
$\mathfrak P_1,\mathfrak P_2$ above the same good prime
$\mathfrak p$ of $K$, the points
\[
 [h]P+T_1,\qquad [h-c]P+T_2
\]
reduce to the identity at $\mathfrak P_1,\mathfrak P_2$,
respectively. Then $[c]Q$ reduces to the identity at
$\mathfrak p$.
In particular, such contacts are possible only at finitely many
primes of $K$, independently of $h,T_1,T_2$.
The two places need not be distinct, unramified, or split.
\end{theorem}
\begin{proof}
Multiplication by $t$ commutes with reduction and kills each
$T_i$. Thus $[h]Q$ reduces to the identity at
$\mathfrak P_1$, and $[h-c]Q$ does so at $\mathfrak P_2$.
This uses only that multiplication is a group-scheme morphism
on a good-reduction model. It does not divide by $t$ and
remains valid when the residue characteristic divides $t$.
Both points belong to $E(K)$. Extension of a residue field is
injective on the points of the reduced curve, so their
reductions already vanish over the residue field of
$\mathfrak p$. Subtracting gives
$[c]\widetilde Q_{\mathfrak p}=O$.

The point $[c]Q$ is a fixed nonidentity point. Choose a fixed
integral Weierstrass model and discard the finitely many primes
dividing its discriminant. At every remaining prime it reduces
to the identity exactly when its abscissa has negative valuation.
Only finitely many prime ideals divide the denominator of this
fixed abscissa. Restoring the discarded finite set proves the
finiteness assertion, including good primes at which the chosen
global equation was nonminimal.
\end{proof}

An equivalent order formulation is useful when the prime splits.
If $m$ is the order of $\widetilde Q_{\mathfrak p}$, the two
contacts imply $m\mid h$ and $m\mid h-c$, hence $m\mid c$.
For adjacent indices $c=1$, they force
$\widetilde Q_{\mathfrak p}=O$.

\begin{corollary}\label{cor:menu}
For any fixed finite set $C\subset\Z\setminus\{0\}$, there is
a finite set of primes outside which none of the contact pairs
in Theorem~\ref{thm:reduction} is possible for any $c\in C$.
\end{corollary}
\begin{proof}
Choose a fixed integral Weierstrass model. Take the union of
the finite denominator supports of $x([c]Q)$ for $c\in C$,
together with the primes dividing the discriminant of that model.
\end{proof}

\section{Rank-stable curves and the conjugate detector}

\begin{corollary}\label{cor:rankstable}
Suppose $E/K$ has
\[
 0<\operatorname{rank}E(K)=\operatorname{rank}E(L).
\]
Then every infinite-order $P\in E(L)$ satisfies the descent
hypothesis of Theorem~\ref{thm:reduction}. If the common rank is
one and $P$ generates the free part of $E(L)$, the obstruction
covers the entire point domain modulo rational torsion.
\end{corollary}
\begin{proof}
The Mordell--Weil groups are finitely generated. Their quotient
$E(L)/E(K)$ has rank zero, hence is finite. A positive integer
annihilating that quotient carries every $P$ into $E(K)$.
In rank one every point is an integer multiple of a free
generator plus torsion.
\end{proof}

For example, let $L/K$ be quadratic and let $\sigma$ be its
nontrivial automorphism. A conjugate adjacent-index detector
requires, at a chosen place $\mathfrak P$ of $L$, simultaneous
positive valuations of the reciprocal abscissas
\[
 \frac1{x(H)},\qquad
 \sigma\left(\frac1{x(H-P)}\right),\qquad H=[h]P,
\]
after excluding the fixed coefficient and bad-reduction primes.
The first condition says that $H$ reduces to the identity at
$\mathfrak P$; the second says that $H-P$ reduces to the identity
at $\sigma^{-1}\mathfrak P$. Theorem~\ref{thm:reduction} with
$c=1$ therefore excludes the pair outside the finite support of
$x(Q)$. A balance transformation preserving positivity of these
two original tests does not alter this conclusion.

Zywina's theorem~\cite[Theorem 1.2]{Zywina} supplies rank-one
curves over $K$ whose rank remains one over a prescribed quadratic
extension $L/K$. Such curves have useful global height properties,
but they satisfy the hypotheses of Corollary~\ref{cor:rankstable}.
They cannot supply arbitrarily large primes for this particular
conjugate adjacent-index preparation. The obstruction does not
concern curves over $L$ that are not defined over $K$, nor does it
exclude a different local comparison construction.

No assertion here concerns the existence of a Diophantine
definition of $\mathcal O_K$ in $K$. The theorem isolates one
failed substitution in a proposed proof of that assertion.

\begin{thebibliography}{9}
\bibitem{Zywina}
D. Zywina, \emph{Rank one elliptic curves and rank stability},
preprint, 2025, arXiv:2505.16960v1, Theorem 1.2.
\url{https://arxiv.org/abs/2505.16960}.
\bibitem{Silverman}
J. H. Silverman, \emph{The Arithmetic of Elliptic Curves},
second edition, Graduate Texts in Mathematics 106,
Springer, 2009; Chapters VII and VIII.
\end{thebibliography}
\end{document}
